\chapter{WordPress}

\section{Introduzione}

\textbf{WordPress} è un CMS (Content Management System) nato come strumento per il blogging, oggi esteso ad ogni tipo di utilizzo. È un software Open Source ed è una delle piattaforme più utilizzate al mondo per la creazione di siti web.

\section{Caratteristiche principali}

Le caratteristiche principali di WordPress sono:

\begin{itemize}
    \item \textbf{Temi}: controllano l'aspetto grafico e funzionale del sito.
    \item \textbf{Plugin}: estendono le funzionalità di WordPress.
    \item \textbf{Ottimizzato per i motori di ricerca} (SEO-friendly).
    \item \textbf{Open source}: il codice sorgente è liberamente disponibile.
\end{itemize}

\section{Funzioni base}

\subsection{Gestione utenti e ruoli}

WordPress gestisce diversi ruoli utente, ciascuno con permessi specifici:

\begin{itemize}
    \item \textbf{Super Admin}: disponibile solo in WordPress Multisite.
    \item \textbf{Administrator}: ha pieno controllo sul sito.
    \item \textbf{Editor}: controllo completo sui contenuti (propri e di altri utenti).
    \item \textbf{Author}: può creare, modificare e pubblicare i propri contenuti.
    \item \textbf{Contributor}: può creare contenuti e modificare i propri, ma non può pubblicarli.
    \item \textbf{Subscriber}: non può creare contenuti.
\end{itemize}

\subsection{Gestione documenti multimediali}

WordPress include un sistema integrato per la gestione di documenti multimediali (immagini, video, audio, PDF, ecc.) tramite la \emph{Media Library}.

\subsection{Sistema di commenti integrato}

WordPress dispone di un sistema di commenti integrato che permette ai visitatori di interagire con i contenuti pubblicati.

\section{Contenuti su WordPress}

WordPress organizza i contenuti in diversi tipi:

\begin{itemize}
    \item \textbf{Post}: articoli del blog.
    \item \textbf{Page}: pagine statiche.
    \item \textbf{Attachment}: documenti multimediali allegati.
    \item \textbf{Revision}: versioni precedenti dei post.
    \item \textbf{Navigation menu}: menu di navigazione.
    \item \textbf{Custom post}: tipi di contenuto personalizzati definiti dall'utente.
\end{itemize}

WordPress fornisce anche una \textbf{REST API} (\url{https://developer.wordpress.org/rest-api/}) che permette a sviluppatori esterni di interagire programmaticamente con il CMS.

\section{Configurazioni: wp-config.php}

Alcune configurazioni principali di WordPress si trovano nel file \texttt{wp-config.php}. Esempi:

\begin{lstlisting}[language=PHP]
define('WP_DEBUG', true);
define('WP_CACHE', true);
\end{lstlisting}

\begin{itemize}
    \item \texttt{WP\_DEBUG}: attiva la modalità di debug per lo sviluppo.
    \item \texttt{WP\_CACHE}: abilita il sistema di caching.
\end{itemize}

\section{Bacheca di WordPress}

La \textbf{bacheca} (o \emph{dashboard}) di WordPress è l'interfaccia di amministrazione del sito. Attraverso la bacheca è possibile gestire:

\begin{itemize}
    \item Post e pagine.
    \item Media.
    \item Commenti.
    \item Temi e plugin.
    \item Utenti.
    \item Impostazioni generali del sito.
\end{itemize}

\section{Organizzazione dei contenuti}

\subsection{I Post}

I \textbf{Post} in WordPress hanno due accezioni:

\begin{itemize}
    \item Sono \textbf{articoli} pubblicati sul sito.
    \item Dal punto di vista interno, sono \textbf{tutti i contenuti} memorizzati nella tabella \texttt{wp\_posts} del database. Il campo \texttt{post\_type} distingue:
    \begin{itemize}
        \item \textbf{Post}: articoli standard.
        \item \textbf{Pagine}: contenuti statici.
        \item \textbf{Post personalizzati}: ad esempio eventi di un calendario, schede dei prodotti, ecc.
        \item \textbf{Revisioni}: versioni precedenti dei post (da salvataggi bozze).
        \item \textbf{Elementi dei menu di navigazione} (link).
    \end{itemize}
\end{itemize}

\subsection{Stato dei post}

Ogni post può trovarsi in uno dei seguenti stati:

\begin{center}
\begin{tabular}{|l|p{9cm}|}
\hline
\textbf{Stato} & \textbf{Descrizione} \\
\hline
Publish & Post pubblicato e visibile pubblicamente. \\
\hline
Future & Post programmato (in attesa di pubblicazione). \\
\hline
Draft & Post in bozza e visibile da utenti con ruolo appropriato (autore e gli utenti Editor o superiori). \\
\hline
Pending & In attesa di revisione da parte di Editor o superiori. \\
\hline
Private & Visibile solo dagli Administrator. \\
\hline
Trash & Cestinato. \\
\hline
Auto-Draft & Salvataggio automatico (revisione). \\
\hline
Inherit & Post che ereditano lo stato da un post genitore (ad es.\ Attachment o Revision rispetto a un Post). \\
\hline
\end{tabular}
\end{center}

\subsection{Permessi sui post}

I permessi variano in base al ruolo dell'utente:

\begin{center}
\begin{tabular}{|l|p{9cm}|}
\hline
\textbf{Ruolo} & \textbf{Capacità} \\
\hline
Subscriber & Non può creare contenuti. \\
\hline
Contributor & Può creare contenuti e modificare i propri, non può pubblicare. \\
\hline
Author & Può creare contenuti, modificare e pubblicare i propri. \\
\hline
Editor & Ha il controllo completo sui contenuti propri e di altri utenti; può pubblicare contenuti in attesa di revisione creati dai Contributor. \\
\hline
\end{tabular}
\end{center}

\section{Tassonomie}

Una \textbf{tassonomia} è un insieme di termini legati da un'analogia semantica e/o gerarchica. Le tassonomie:

\begin{itemize}
    \item Permettono di assegnare ai contenuti \textbf{parole chiave con valore semantico}; tali termini consentono all'utente di comprendere immediatamente il contenuto di un articolo e di effettuare ricerche per argomenti (\emph{organizzazione semantica}).
    \item Permettono di creare un \textbf{sistema di navigazione} facilmente comprensibile da parte dei visitatori (\emph{organizzazione gerarchica}).
\end{itemize}

\subsection{Categorie}

Le \textbf{Categorie} consentono di raggruppare logicamente e gerarchicamente i contenuti del blog. Costituiscono una \emph{tassonomia gerarchica}.

\subsection{Tag}

I \textbf{Tag} etichettano i contenuti in modo tale che l'utente comprenda immediatamente quale sia l'argomento di un articolo. Costituiscono una \emph{tassonomia semantica}.

\section{Pagine}

Le \textbf{Pagine}:

\begin{itemize}
    \item Sono \textbf{contenuti statici} (non articoli del blog).
    \item Possono avere una \textbf{struttura gerarchica}: può essere definito il genitore in caso di nuova pagina, creando relazioni padre-figlio.
\end{itemize}

Sono adatte per contenuti come ``Chi siamo'', ``Contatti'', ``Termini di servizio'', ecc.

\section{Plugin}

\subsection{Funzionalità core e Plugin}

WordPress offre delle funzionalità \emph{core}, ma i \textbf{Plugin} permettono di aggiungere funzionalità ulteriori.

Un plugin può essere costituito da:

\begin{itemize}
    \item Un solo file principale.
    \item Una collezione di file \texttt{.php}, \texttt{.css}, \texttt{.js}, oltre a eventuali grafiche.
\end{itemize}

\subsection{Widget}

I \textbf{Widget} sono blocchi di codice HTML che permettono di aggiungere contenuto o funzionalità alle \emph{sidebar} di WordPress. WordPress dispone di diversi widget predefiniti per:

\begin{itemize}
    \item Categorie.
    \item Tag cloud.
    \item Ricerca.
    \item Calendario delle pubblicazioni.
    \item Menu personalizzati.
    \item Ecc.
\end{itemize}

\subsection{Shortcode}

Gli \textbf{Shortcode} sono particolari tag che vengono inseriti all'interno dei contenuti del sito ed eseguono delle macro, producendo contenuto HTML.

Esempio:

\begin{verbatim}
[gallery ids="123,124,125,126" size="medium"]
\end{verbatim}

Questo shortcode, inserito in un post o in una pagina, genera una galleria di immagini con gli ID specificati.

\section{Temi}

I \textbf{Temi} sono le estensioni di WordPress che controllano il modo in cui un sito si presenta ai visitatori e agli utenti.

Non si tratta solo della definizione delle caratteristiche grafiche, ma del \textbf{controllo generale di tutti gli aspetti che riguardano la presentazione}, anche sotto il profilo funzionale.

Un tema determina:

\begin{itemize}
    \item Il layout del sito.
    \item La tipografia e i colori.
    \item La struttura delle pagine e dei post.
    \item Le aree widget disponibili.
    \item I template per i vari tipi di contenuto.
\end{itemize}
